\documentclass[11pt,a4paper]{article}
\usepackage{amsmath,amssymb}
\usepackage{listings}
\usepackage{hyperref}
\usepackage[margin=1in]{geometry}
\usepackage{fancyhdr}

\pagestyle{fancy}
\fancyhf{}
\fancyhead[L]{Module 03: Control Flow}
\fancyhead[R]{C Programming Course}
\fancyfoot[C]{\thepage}

\lstset{
  language=C,
  basicstyle=\ttfamily\small,
  keywordstyle=\bfseries,
  commentstyle=\itshape,
  numbers=left,
  numberstyle=\tiny,
  frame=single,
  breaklines=true,
  tabsize=4
}

\title{Module 03: Control Flow}
\author{C Programming Course}
\date{}

\begin{document}
\maketitle
\thispagestyle{fancy}

\section{Conditional Statements}
\subsection{if, else, else-if}
\begin{lstlisting}
#include <stdio.h>

int main(void) {
    int score = 75;
    if (score >= 90) {
        printf("Grade: A\n");
    } else if (score >= 80) {
        printf("Grade: B\n");
    } else if (score >= 70) {
        printf("Grade: C\n");
    } else {
        printf("Grade: F\n");
    }
    return 0;
}
\end{lstlisting}

\subsection{Switch Statement}
\begin{lstlisting}
#include <stdio.h>

int main(void) {
    int day = 3;
    switch (day) {
        case 1: printf("Monday\n");    break;
        case 2: printf("Tuesday\n");   break;
        case 3: printf("Wednesday\n"); break;
        default: printf("Other\n");    break;
    }
    return 0;
}
\end{lstlisting}

\section{Loops}
\subsection{for Loop}
\begin{lstlisting}
for (int i = 0; i < 10; i++) {
    printf("%d ", i);
}
\end{lstlisting}

\subsection{while Loop}
\begin{lstlisting}
int n = 5;
while (n > 0) {
    printf("%d ", n);
    n--;
}
\end{lstlisting}

\subsection{do-while Loop}
Executes the body at least once before checking the condition.
\begin{lstlisting}
int x = 0;
do {
    printf("%d ", x);
    x++;
} while (x < 3);
\end{lstlisting}

\section{Break and Continue}
\begin{itemize}
  \item \texttt{break} -- exits the innermost loop or switch immediately.
  \item \texttt{continue} -- skips the rest of the current iteration.
\end{itemize}

\section{Nested Loops and Loop Performance}
\begin{lstlisting}
#include <stdio.h>

int main(void) {
    for (int i = 1; i <= 3; i++) {
        for (int j = 1; j <= 3; j++) {
            printf("(%d,%d) ", i, j);
        }
        printf("\n");
    }
    return 0;
}
\end{lstlisting}

Nested loops multiply iteration counts: an $O(n)$ loop inside another $O(n)$
loop yields $O(n^2)$ total iterations.

\end{document}
